\begin{abstract}

Rotation invariance is essential for precise object level segmentation in UAV aerial imagery, where targets can have arbitrary orientations and exhibit fine scale details. Conventional segmentation architectures like UNet rely on convolution operators that are not rotation invariant, leading to degraded segmentation accuracy across varying viewpoints. Rotation invariance can be achieved by expanding the filter bank across multiple orientations; however, this significantly increases computational cost and memory requirement.
In this paper, we introduce a GPU optimized rotation invariant convolution framework that eliminates the traditional data lowering (im2col) step required for matrix multiplication based convolution. By exploiting structured data sharing among symmetrically rotated filters, our method achieves multi orientation convolution with greatly reduced memory requirements and computational redundancy. We further generalize the approach to accelerate convolution with arbitrary (non symmetric) rotation angles.
Integrated into a UNet segmentation model, the framework yields up to a 5.7\% improvement in accuracy over the non rotation aware baseline. Across extensive benchmarks, the proposed convolution achieves 20–57\% faster training and 15–45\% lower energy consumption than cuDNN, while maintaining accuracy comparable to state-of-the-art rotation invariant methods. Because the scatter based operator greatly reduces intermediate feature dimensionality, the efficiency of our design also enables practical sixteen orientation convolution and pooling, yielding further accuracy gains that are infeasible for conventional rotation invariant implementations. These results demonstrate that the proposed method provides an effective and highly efficient alternative to existing rotation invariant convolution frameworks.


%Across extensive benchmarks, the proposed convolution achieves 20–55\% faster training and 15–45\% lower energy consumption than cuDNN, while maintaining accuracy comparable to state-of-the-art rotation invariant methods. In the eight-orientation setting, our approach achieves up to 45\% speedup and 41\% energy savings on 256×256 inputs, and 32\% speedup and 23\% lower energy usage on 1024×1024 inputs. Integrated into a U-Net segmentation model, the framework yields up to 6\% improvement in accuracy over the non-rotation-aware baseline. These results demonstrate that the proposed method provides an effective and highly efficient alternative to existing rotation invariant CNN frameworks.
%
%The efficiency of our scatter-based design also enables practical sixteen-orientation convolution and pooling, yielding further accuracy gains that are infeasible for conventional rotation-invariant implementations.

%In benchmarks against cuDNN-based convolution implementations, the proposed method delivers higher throughput and lower energy consumption across most evaluated workloads. In our tests, it achieved $\sim$1.8× speedup and $\sim$34\% reduction in energy use. Furthermore, when compared with the baseline network without rotation handling, the proposed framework achieved $\sim$2\% improvement in segmentation accuracy.

%Rotation invariance is essential for precise, object-level segmentation in UAV aerial imagery, where targets can have arbitrary orientations and exhibit fine-scale details. Conventional segmentation architectures like U-Net rely on convolution operators that are not rotation-invariant. Implementing rotation invariance by expanding the filter bank across multiple orientations significantly increases computational cost and memory traffic.
%We introduce a GPU-optimized rotation-invariant convolution framework that eliminates the traditional data-lowering (im2col) step required for matrix-multiplication-based convolution. By exploiting structured data sharing among symmetrically rotated filters, our method achieves multi-orientation convolution with significantly reduced memory traffic and computational redundancy.
%We further generalize the approach to accelerate convolution with arbitrary (non‑symmetric) rotation angles. In benchmarks against cuDNN based implementation, our implementation delivers higher throughput and lower energy consumption across most evaluated workloads. In our tests, it achieved $\sim$1.8$\times$ speedup and $\sim$34$\%$ reduction in energy use.

%Convolution is a foundational operation in convolution-based deep learning. This paper presents a new approach for implementing dense convolution by leveraging matrix multiplication on the GPU. Unlike existing methods, the proposed approach eliminates the data-lowering (im2col) step that is typically required when applying matrix multiplication to implement convolution. Standard convolution is not inherently rotation-invariant. Rotation invariance can be achieved by rotating each filter so the network learns features at multiple orientations, but this increases training time because the number of convolutions grows with the number of rotations. Our method accelerates rotation-invariant convolution for symmetric rotations by exploiting the operation’s inherent data-sharing properties, and we further extend it to handle arbitrary rotations. We benchmark our proposed convolution against cuDNN based convolution implementation. The experiments show that our method delivers higher throughput and lower energy consumption in most test cases.

%Convolution is central to deep learning, yet standard implementations are not rotation‑invariant and become computationally expensive when naively extended to multiple orientations. We present a GPU‑optimized rotation‑invariant convolution framework that avoids the conventional data‑lowering (im2col) step typically required for matrix‑multiplication‑based convolution. By exploiting the inherent data‑sharing structure among symmetrically rotated filters, our method performs multi‑orientation convolution with substantially reduced memory traffic and arithmetic redundancy; we further generalize the approach to arbitrary (non‑symmetric) rotations. The same kernel also accelerates conventional dense convolution in the single‑orientation case. In benchmarks against cuDNN, our implementation delivers higher throughput and lower energy consumption across most tested workloads.
\end{abstract}

\begin{IEEEkeywords}
Dense convolution, Scatter operation, Rotation invariant, Acceleration, GPU
\end{IEEEkeywords}





